\section{2024-11-11 - Das Versagen der UNO}

\subsection{Die zentrale Aussage des Autors}

Bei dem Kommentar \quotation{Das klägliche Versagen der UNO}, verfasst von
Andreas Zumach, veröfentlicht am 24.09.2016 im "Deutschlandfunk", befasst sich
mit der Frage, ob die UNO noch handlhungsfähig bzw. legitim ist.

Zumach beginnt mit dem nach ihm kläglichen Versagen der UNO (vgl. Z. 11) die
eskalierung des Syrienkrieges zu unterbinden oder auch nur umanitäre Hilfe für
die geflüchteten Menschen zu gewährleisten.

Zentrale Aussagen
\startitemize[packed]
\item Die UNO hat dabei versagt:
\startitemize[packed]
\item den \quotation{eskalierenden Krieg in Syrien endlich zu beenden} (Z. 8)
\item die geflüchteten Menschen \quotation{humanitär zu versorgen} (Z. 10)
\stopitemize
\item im Detail das Versagen der USA und Russland, die nach 
\startitemize[packed]
\item Beide wollen den Islamischen Staat bekämpfen
\item USA eher auf Seite der Opposition
\item gegenseitig erkennen die ihre verbündeten nicht an
\stopitemize
\stopitemize


\startframedtipp
Hintergrundinformation in aufgabe 2 (Stellungnahem) mit aufnehmen, damit er sieht, was wir wissen

\startitemize[packed]
\item Gut überhaupt internationales zu habe
\item andere gute aspekte wie unicef etc
\item Leistet auch woander humatiäre hilfe, wenn auch nicht genug (Besser als nichts)
\stopitemize
\stopframedtipp
